% !TEX program = xelatex
\documentclass[11pt,a4paper]{article}

\usepackage[UTF8]{ctex}
\usepackage[margin=2.2cm]{geometry}
\usepackage{amsmath, amssymb, amsthm}
\usepackage{listings}
\usepackage{xcolor}
\usepackage{tikz}
\usepackage{booktabs}
\usepackage{multirow}
\usepackage{hyperref}
\usetikzlibrary{positioning, arrows.meta, shapes.geometric, fit, calc, decorations.pathreplacing}

\hypersetup{
  colorlinks=true,
  linkcolor=blue!60!black,
  urlcolor=teal!70!black,
}

\lstdefinelanguage{Rust}{
  morekeywords={fn,let,mut,pub,struct,impl,use,mod,self,Self,Option,Some,None,return,if,else,for,while,loop,match,break,continue,as,where,trait,type,enum,move,ref,const,static,unsafe,extern,crate,super,dyn,async,await,box,in},
  sensitive=true,
  morecomment=[l]{//},
  morecomment=[s]{/*}{*/},
  morestring=[b]",
}

\lstset{
  language=Rust,
  basicstyle=\ttfamily\footnotesize,
  keywordstyle=\color{blue!60!black}\bfseries,
  commentstyle=\color{green!40!black}\itshape,
  stringstyle=\color{red!60!black},
  numbers=left,
  numberstyle=\tiny\color{gray},
  numbersep=8pt,
  frame=single,
  framerule=0.4pt,
  rulecolor=\color{gray!40},
  breaklines=true,
  breakatwhitespace=true,
  showstringspaces=false,
  tabsize=4,
  columns=fullflexible,
  keepspaces=true,
}

\title{\textbf{halfedge\,: 布尔运算模块设计文档} \\ \large \texttt{src/boolean.rs}}
\author{halfedge 库}
\date{2026-07-06}

\begin{document}
\maketitle
\tableofcontents

\section{模块定位}

\texttt{boolean.rs} 在两个闭合流形三角网格 $A$、$B$ 之上执行\textbf{集合运算}
（并、交、差、对称差），返回新的 \texttt{MeshStorage}。

\textbf{算法路线}：\textbf{Corefinement（共精化）}。对每个三角形找边-三角形交点，
在交点处精确切割三角形为子三角形，再对子三角形做内外分类后按布尔规则保留/丢弃。
与旧射线法相比，corefinement 对跨越边界的三角形进行精确几何切割，
使 $A - A$ \textbf{严格为空}（0 面），并提升部分重叠场景的几何精度。

\subsection{API 总览}

\begin{center}
\renewcommand{\arraystretch}{1.18}
\resizebox{\textwidth}{!}{%
\begin{tabular}{@{}lll@{}}
  \toprule
  \textbf{类别} & \textbf{API} & \textbf{语义} \\
  \midrule
  枚举   & \texttt{BoolOp}                            & \texttt{Union}/\texttt{Intersection}/\texttt{Difference}/\texttt{SymmetricDifference} \\
  核心   & \texttt{boolean\_operation(a, b, op)}      & 按 \texttt{op} 分派，返回新 \texttt{MeshStorage} \\
  \midrule
  \multirow{4}{*}{快捷}
         & \texttt{boolean\_union(a, b)}              & $A \cup B$ \\
         & \texttt{boolean\_intersection(a, b)}       & $A \cap B$ \\
         & \texttt{boolean\_difference(a, b)}         & $A - B$ \\
         & \texttt{boolean\_symmetric\_difference(a, b)} & $A \triangle B$ \\
  \bottomrule
\end{tabular}%
}
\end{center}

\section{BoolOp 分类规则}

\texttt{classify(op, from\_a, inside\_other)} 根据子三角形来源与在另一网格
内外的状态决定保留：

\begin{center}
\renewcommand{\arraystretch}{1.18}
\begin{tabular}{@{}lll@{}}
  \toprule
  \textbf{op} & \textbf{A 面 (from\_a=true)} & \textbf{B 面 (from\_a=false)} \\
  \midrule
  \texttt{Union}              & \texttt{!inside\_other} & \texttt{!inside\_other} \\
  \texttt{Intersection}       & \texttt{inside\_other}  & \texttt{inside\_other}  \\
  \texttt{Difference}         & \texttt{!inside\_other} & \texttt{false}（永不保留） \\
  \texttt{SymmetricDifference}& \texttt{!inside\_other} & \texttt{!inside\_other} \\
  \bottomrule
\end{tabular}
\end{center}

\textbf{语义}：保留\textbf{位于结果边界上}的子三角形。例如并集保留 $A$ 表面在
$B$ 外部部分 + $B$ 表面在 $A$ 外部部分；差集仅保留 $A$ 表面在 $B$ 外部部分。
这与 $A - A = \emptyset$ 的几何直觉一致：$A$ 表面在 $A$ 内部 $\Rightarrow$
\texttt{inside\_other=true} $\Rightarrow$ \texttt{!inside\_other=false} $\Rightarrow$ 全部丢弃。

\section{核心算法：Corefinement 四阶段}

\subsection{阶段 1：边-三角形求交}

对 $A$ 的每个三角形 $T$，遍历其三条边 $e_k$ 与 $B$ 的所有三角形 $T'$，
检测交点。分两种情形：

\subsubsection{非共面情形：\texttt{segment\_triangle\_intersection}}

设线段端点 $p_0, p_1$，三角形顶点 $v_0, v_1, v_2$。
用 Shewchuk 鲁棒 \texttt{orient3d} 谓词计算两端符号：
\[
  s_0 = \mathrm{orient3d}(v_0, v_1, v_2, p_0),\quad
  s_1 = \mathrm{orient3d}(v_0, v_1, v_2, p_1).
\]
\begin{itemize}
  \item 同号 $\Rightarrow$ 无交点；
  \item 共面（$s_0 = s_1 = 0$）$\Rightarrow$ 转共面处理；
  \item 端点在平面上（$s_0 = 0$ 或 $s_1 = 0$）$\Rightarrow$ 跳过，避免重复顶点；
  \item 异号 $\Rightarrow$ 求交点参数：
  \[
    t = \frac{s_0}{s_0 - s_1} \in (10^{-9},\; 1 - 10^{-9}).
  \]
\end{itemize}
求得交点 $P = p_0 + t(p_1 - p_0)$ 后，投影到三角形法向主轴的正交 2D 平面，
用 \texttt{point\_in\_triangle\_2d} 判定是否在三角形内。

\subsubsection{共面情形：\texttt{coplanar\_segment\_segment\_3d}}

当线段与三角形共面时，对三角形的每条边做 2D 线段-线段求交。投影到两条
边叉积的主轴正交平面，用参数方程求解：
\[
  t_1 = \frac{(p_3 - p_1) \times d_2}{d_1 \times d_2},\quad
  t_2 = \frac{(p_3 - p_1) \times d_1}{d_1 \times d_2},
\]
其中 $d_1 = p_2 - p_1$，$d_2 = p_4 - p_3$。要求 $t_1, t_2 \in (10^{-9}, 1-10^{-9})$，
端点相交被过滤。

\subsubsection{完整共面重叠：\texttt{coplanar\_triangle\_overlap}}

当三角形 $A$ 的\textbf{三个顶点}均在 $B$ 的平面上时（\texttt{orient3d} 均为 0），
两面完整共面。此时边-边求交无法捕捉\textbf{面内包含}情形（如 $B$ 整体落在
$A$ 内部，无边交叉），需用 Sutherland-Hodgman 多边形裁剪计算重叠区域。

\textbf{算法}：
\begin{enumerate}
  \item 投影到 2D（丢弃法向主轴，与 \texttt{coplanar\_segment\_segment\_3d} 一致）；
  \item 检测 $B$ 的 2D 朝向，CW 时反转边方向确保 CCW；
  \item 逐边将 $A$ 裁剪到 $B$ 内部（Sutherland-Hodgman）：
    \begin{itemize}
      \item 对 $B$ 的每条边 $(e_s, e_e)$，保留 $A$ 中在边左侧（CCW 内部）的顶点，
            计算进出交点；
      \item 3 次裁剪后得到 $A \cap B$ 的凸多边形；
    \end{itemize}
  \item 利用平面方程 $\vec{n} \cdot \vec{x} + d = 0$ 将 2D 顶点反解回 3D。
\end{enumerate}

\textbf{约束点提取}：重叠多边形的顶点分为三类：
\begin{itemize}
  \item $A$ 的顶点在 $B$ 内部 $\Rightarrow$ 已是 $A$ 顶点，无需额外处理；
  \item $B$ 的顶点在 $A$ 内部 $\Rightarrow$ \textbf{内部约束点}，加入
        \texttt{interior\_points} 列表；
  \item 边-边交点 $\Rightarrow$ 已在 $A$ 的边上，加入对应边的交点列表。
\end{itemize}

\textbf{内部约束点的分裂}：在阶段 2 的耳裁剪三角化之后，对每个内部约束点 $p$，
定位包含 $p$ 的子三角形 $T' = (v_0, v_1, v_2)$，做\textbf{扇形分裂}：
\[
  T' \to (v_0, v_1, p),\;(v_1, v_2, p),\;(v_2, v_0, p).
\]
多个内部约束点逐个处理，每次在当前子三角形集合中定位并分裂。

\subsection{阶段 2：三角形分裂}

将每条边上的交点按参数 $t$ 排序后插入三角形边界，形成多边形：
\[
  \text{polygon} = [v_0,\; \text{intersections on } e_0,\; v_1,\;
  \text{intersections on } e_1,\; v_2,\; \text{intersections on } e_2].
\]
用 \texttt{ear\_clipping\_3d} 对多边形做耳裁剪三角化，得到子三角形列表。
随后对每个\textbf{内部约束点}（来自共面重叠）做扇形分裂（见上节）。
无交点且无内部约束点时返回原三角形。

\subsection{阶段 3：内外分类（多采样 + 内向偏移）}

对每个子三角形 $T'$ 采样 7 个点：
\[
  \mathcal{S} = \{\text{重心}\} \cup \{3 \text{ 顶点}\} \cup \{3 \text{ 边中点}\}.
\]
沿子三角形法向 $\hat{n}$ 的\textbf{内向}偏移 $\varepsilon$：
\[
  \vec{p}_{\text{sample}} = \vec{s} - \varepsilon \hat{n},\quad
  \varepsilon = \frac{10^{-4}}{\|\vec{n}\|}\;\;(\|\vec{n}\| < 10^{-14} \text{ 时回退 } 10^{-4}).
\]
对每个采样点用射线奇偶法判定是否在另一网格内部，\textbf{多数表决}
（$\ge 4/7$）$\Rightarrow$ \texttt{inside\_other=true}。

\textbf{内向偏移的意义}：对共面情形（如 $A - A$），三角形恰在另一网格表面，
单侧采样会产生歧义。内向偏移确保采样点明确落入网格内部，使
$A - A$ 严格为空。

\textbf{共面回退}：当全部 7 个采样点均判定为外部（$\text{inside\_count} = 0$）时，
可能是\textbf{背对背共面}（$A$ 的面法向 $+z$，$B$ 的面法向 $-z$），内向偏移使
采样点离开对方体积。此时用 2D 包含检测回退：若子三角形重心与网格某面共面
（\texttt{orient3d} 均为 0）且在该面的 2D 投影内，则判定为 \texttt{inside}。
此回退确保并集/差集中背对背共面区域被正确丢弃，避免非流形结果。

\subsection{阶段 4：集合运算与重建}

对每个子三角形调用 \texttt{classify(op, from\_a, inside\_other)} 决定保留。
保留的子三角形通过 \texttt{get\_or\_add\_vertex}（量化 $10^{-9}$ 去重）
+ \texttt{add\_triangle} 重建结果网格。

\subsection{整体流程}

\begin{center}
\resizebox{\textwidth}{!}{%
\begin{tikzpicture}[
  node distance=0.4cm,
  proc/.style={draw, rounded corners=2pt, minimum width=12cm, minimum height=0.6cm, align=center, font=\small},
  op/.style={-Stealth, thick},
  startend/.style={draw, rounded corners=10pt, minimum width=2.6cm, minimum height=0.6cm, font=\small, fill=gray!12}
]
  \node[startend] (s) {输入：$A$, $B$, \texttt{op}};
  \node[proc, below=of s, fill=blue!8] (p1) {\textbf{阶段 1}：边-三角形求交；共面时 Sutherland-Hodgman 计算重叠多边形};
  \node[proc, below=of p1, fill=blue!8] (p2) {\textbf{阶段 2}：交点 + 内部约束点 $\Rightarrow$ 耳裁剪 + 扇形分裂三角化};
  \node[proc, below=of p2, fill=blue!8] (p3) {\textbf{阶段 3}：7 采样 + 内向偏移 $\varepsilon$ + 多数表决；共面回退 2D 包含};
  \node[proc, below=of p3, fill=blue!8] (p4) {\textbf{阶段 4}：\texttt{classify(op, from\_a, inside\_other)} 决定保留};
  \node[proc, below=of p4, fill=blue!8] (p5) {对 $B$ 的每个三角形重复阶段 1--4};
  \node[proc, below=of p5, fill=blue!8] (p6) {保留子三角形经量化去重 + \texttt{add\_triangle} 重建};
  \node[startend, below=of p6, fill=green!12] (e) {返回新 \texttt{MeshStorage}};
  \draw[op] (s) -- (p1);
  \draw[op] (p1) -- (p2);
  \draw[op] (p2) -- (p3);
  \draw[op] (p3) -- (p4);
  \draw[op] (p4) -- (p5);
  \draw[op] (p5) -- (p6);
  \draw[op] (p6) -- (e);
\end{tikzpicture}%
}
\end{center}

\subsection{数据丢弃警告}

当 \texttt{add\_triangle} 创建三角形失败时，不再静默丢弃，而是计数并在返回前
通过 \texttt{eprintln!} 输出警告：
\begin{center}
\texttt{[halfedge::boolean] 警告：N 个三角形创建失败（拓扑冲突），已跳过}
\end{center}

\section{辅助函数}

\begin{center}
\renewcommand{\arraystretch}{1.18}
\resizebox{\textwidth}{!}{%
\begin{tabular}{@{}lll@{}}
  \toprule
  \textbf{函数} & \textbf{签名} & \textbf{用途} \\
  \midrule
  \texttt{segment\_triangle\_intersection}
    & \texttt{(p0, p1, v0, v1, v2) -> Option<(V3, f64)>}
    & 非共面线段-三角形求交 \\
  \texttt{point\_in\_triangle\_3d}
    & \texttt{(p, v0, v1, v2) -> bool}
    & 投影 2D 后用 \texttt{point\_in\_triangle\_2d} 判定 \\
  \texttt{coplanar\_segment\_segment\_3d}
    & \texttt{(e1a, e1b, e2a, e2b) -> Option<(V3, f64, f64)>}
    & 共面线段-线段 2D 参数求交 \\
  \texttt{sutherland\_hodgman\_clip}
    & \texttt{(\&[2D], 2D, 2D) -> Vec<[2D]>}
    & Sutherland-Hodgman 多边形裁剪（半平面） \\
  \texttt{coplanar\_triangle\_overlap}
    & \texttt{(a: [V3;3], b: [V3;3]) -> Option<Vec<V3>>}
    & 共面三角形重叠多边形（3D 顶点，CCW） \\
  \texttt{is\_point\_strictly\_inside\_triangle}
    & \texttt{(p, tri) -> bool}
    & 判定点在三角形内部（排除边界） \\
  \texttt{point\_on\_triangle\_edge}
    & \texttt{(p, tri) -> Option<(usize, f64)>}
    & 找点所在边及参数 $t \in (0,1)$ \\
  \texttt{find\_containing\_triangle}
    & \texttt{(tris, p) -> Option<usize>}
    & 在子三角形列表中定位包含 $p$ 的三角形 \\
  \texttt{point\_inside\_mesh}
    & \texttt{(point, mesh) -> bool}
    & 射线奇偶法（$+x$ 方向） \\
  \texttt{is\_triangle\_inside\_mesh}
    & \texttt{(tri, mesh, \&Bvh) -> bool}
    & 7 采样 + 内向偏移 + 多数表决；共面回退 2D 包含 \\
  \texttt{collect\_triangle\_intersections}
    & \texttt{(tri, mesh, \&Bvh) -> ([Vec<EdgeIntersection>; 3], Vec<V3>)}
    & 收集三边交点 + 共面内部约束点；BVH 候选过滤 \\
  \texttt{split\_triangle\_by\_intersections}
    & \texttt{(tri, edge\_int, \&[V3]) -> Vec<[V3; 3]>}
    & 插入交点 + 耳裁剪 + 内部约束点扇形分裂 \\
  \bottomrule
\end{tabular}%
}
\end{center}

\section{顶点去重}

\texttt{get\_or\_add\_vertex} 通过量化网格去重：
\[
  \text{key} = \bigl(\lfloor 10^{9} \cdot x \rceil,\ \lfloor 10^{9} \cdot y \rceil,\ \lfloor 10^{9} \cdot z \rceil\bigr) \in \mathbb{Z}^3.
\]
\texttt{HashMap<[i64;3], VertexId>} 缓存已添加顶点；阈值 $10^{-9}$ 内的位置视为同一点。
这保证子三角形间的共享顶点（含 corefinement 产生的交点）能被识别，
使 \texttt{add\_triangle} 自动配对 twin。

\section{Panic 守卫}

逐条检查可能导致 panic 的路径：

\begin{center}
\renewcommand{\arraystretch}{1.18}
\begin{tabular}{@{}ll@{}}
  \toprule
  \textbf{风险路径} & \textbf{守卫} \\
  \midrule
  空网格输入                & \texttt{face\_count() == 0} 时返回空结果 \\
  退化三角形（共线/重合）   & \texttt{is\_triangle\_degenerate\_3d} 跳过 \\
  子三角形退化（耳裁剪产物）& 同上，二次跳过 \\
  零长度法向                & $\|\vec{n}\| < 10^{-14}$ 时 $\varepsilon$ 回退 $10^{-4}$ \\
  端点重复交点              & $t \in (10^{-9}, 1-10^{-9})$ 过滤 \\
  共面边平行                & $|d_1 \times d_2| < 10^{-14}$ 返回 \texttt{None} \\
  多边形顶点不足            & \texttt{polygon.len() < 3} 返回原三角形 \\
  耳裁剪失败                & \texttt{tri\_indices.is\_empty()} 返回原三角形 \\
  \texttt{partial\_cmp} NaN & \texttt{unwrap\_or(Equal)} \\
  \bottomrule
\end{tabular}
\end{center}

\section{复杂度}

设 $F_a = $ \texttt{mesh\_a.face\_count()}，$F_b = $ \texttt{mesh\_b.face\_count()}，
$K = $ 保留子三角形数。

\paragraph{BVH 加速}
\texttt{boolean\_operation} 在入口对每个网格构建一次 BVH
（参见 \texttt{bvh.rs}），阶段 1 与阶段 3 的内层扫描通过
\texttt{Bvh::faces\_in\_aabb}（候选过滤）与 \texttt{Bvh::ray\_count\_hits}
（射线奇偶计数）加速。下表给出加速后的复杂度，$k$ 为候选面数（远小于 $F$）。

\begin{center}
\renewcommand{\arraystretch}{1.18}
\begin{tabular}{@{}ll@{}}
  \toprule
  \textbf{函数} & \textbf{时间复杂度} \\
  \midrule
  \texttt{segment\_triangle\_intersection} & $O(1)$ \\
  \texttt{coplanar\_segment\_segment\_3d}  & $O(1)$ \\
  \texttt{collect\_triangle\_intersections}& $O(\log F_b + k_b)$（每边 BVH 候选 + 精筛） \\
  \texttt{split\_triangle\_by\_intersections}& $O(m)$（$m$ = 交点数，耳裁剪） \\
  \texttt{is\_triangle\_inside\_mesh}      & $O(7 \cdot (\log F_b + k_b))$（7 采样 $\times$ BVH 射线） \\
  \texttt{classify}                        & $O(1)$ \\
  \texttt{get\_or\_add\_vertex}            & $O(1)$ 均摊 \\
  \texttt{Bvh::build}                      & $O(F \log F)$（入口一次性） \\
  \texttt{boolean\_operation}              & $O((F_a + F_b)\log(F_a + F_b) + (F_a + F_b)\bar{k}\bar{m})$ \\
  \bottomrule
\end{tabular}
\end{center}

\textbf{对比}：未加速版本 \texttt{collect\_triangle\_intersections} 与
\texttt{is\_triangle\_inside\_mesh} 均为 $O(F_b)$ 内层扫描，
整体 \texttt{boolean\_operation} 退化为 $O(F_a \cdot F_b \cdot \bar{m})$。
对千面量级网格（$F_a = F_b = 1000$），未加速 $\sim 10^9$ 次操作，
BVH 加速后 $\sim 10^5$ 次操作，约 $10^4\times$ 加速。

\section{限制}

\begin{itemize}
  \item 仅支持\textbf{闭合流形三角网格}；非闭合网格的内外判定无定义；
  \item \texttt{add\_triangle} 的 \texttt{Result} 被显式忽略（\texttt{let \_ = ...}），
        拓扑错误被静默吞掉；
  \item 共面回退对交集运算可能产生非流形结果：两个仅共面接触的网格做交集时，
        两侧的共面面片均被保留（各判定为 inside），产生重复面。
\end{itemize}

\section{测试覆盖}

\begin{center}
\renewcommand{\arraystretch}{1.18}
\resizebox{\textwidth}{!}{%
\begin{tabular}{@{}ll@{}}
  \toprule
  \textbf{测试名} & \textbf{验证点} \\
  \midrule
  \texttt{union\_disjoint\_cubes}                       & 不相交两立方体并集：$F = 24$（无丢弃） \\
  \texttt{union\_overlapping\_cubes}                    & 重叠两立方体并集：$4 \le F \le 24$（内部面被切割丢弃） \\
  \texttt{intersection\_overlapping\_cubes}             & 重叠立方体交集非空（$V > 0$） \\
  \texttt{intersection\_disjoint\_cubes}                & 不相交立方体交集为空（$F = 0$） \\
  \texttt{difference\_self\_is\_empty}                  & $A - A$ \textbf{严格}空（$F = 0$，corefinement 精确切割） \\
  \texttt{symmetric\_difference\_disjoint\_equals\_union} & 不相交时对称差面数 $=$ 并集面数 \\
  \texttt{intersection\_contains\_cube}                 & 两相同立方体交集 $\ge 6$ 面 \\
  \midrule
  \texttt{coplanar\_overlap\_containment}               & 大三角包含小三角（共面）：\texttt{Some} 重叠多边形面积正确 \\
  \texttt{coplanar\_overlap\_partial}                   & 部分共面重叠（共享一边）：\texttt{Some}，面积 $> 0$ \\
  \texttt{coplanar\_overlap\_disjoint}                  & 共面但不相交：\texttt{None} \\
  \texttt{coplanar\_overlap\_non\_coplanar}             & 非共面：\texttt{None} \\
  \texttt{coplanar\_difference\_small\_on\_large}       & 小立方体贴大立方体顶面做差集：拓扑有效，面数减少 \\
  \texttt{coplanar\_union\_small\_on\_large}            & 小立方体贴大立方体顶面做并集：拓扑有效（背对背共面被丢弃） \\
  \texttt{coplanar\_intersection\_small\_on\_large}     & 小立方体贴大立方体顶面做交集：结果非空（仅共面接触） \\
  \midrule
  \multicolumn{2}{@{}l}{\textit{空网格测试}} \\
  \texttt{union\_empty\_meshes\_returns\_empty}        & 双空网格并集：$F = 0$（早退路径） \\
  \texttt{intersection\_empty\_meshes\_returns\_empty} & 双空网格交集：$F = 0$ \\
  \texttt{difference\_empty\_meshes\_returns\_empty}   & 双空网格差集：$F = 0$ \\
  \texttt{union\_empty\_and\_cube\_returns\_cube\_face\_count} & 空 $\cup$ cube：$F \ge 12$（cube 表面在外部全保留） \\
  \texttt{difference\_cube\_and\_empty\_returns\_cube} & cube $-$ 空：$F \ge 12$（Difference 仅保留 A 面在外部） \\
  \texttt{intersection\_cube\_and\_empty\_returns\_empty} & cube $\cap$ 空：$F = 0$（空网格内外判定恒为 false） \\
  \midrule
  \multicolumn{2}{@{}l}{\textit{完全重合测试}} \\
  \texttt{union\_identical\_cubes}                     & 两相同 cube 并集：拓扑有效（内向偏移使表面互判 inside，结果可能为空） \\
  \texttt{intersection\_identical\_cubes}              & 两相同 cube 交集：$F \ge 6$（已知非流形局限，仅校验面数） \\
  \texttt{symmetric\_difference\_identical\_cubes\_is\_empty} & 两相同 cube 对称差：$F = 0$ \\
  \midrule
  \multicolumn{2}{@{}l}{\textit{包含关系测试}} \\
  \texttt{difference\_large\_minus\_small}             & 大 $[-1,1]^3$ 减小 cube（边长 0.5）：拓扑有效 \\
  \texttt{intersection\_large\_and\_small\_returns\_smallish} & 大 $\cap$ 小：$F > 0$，拓扑有效（结果为小 cube 表面） \\
  \midrule
  \multicolumn{2}{@{}l}{\textit{非轴对齐测试}} \\
  \texttt{union\_rotated\_cubes}                       & 轴对齐 cube 与绕 $Z$ 旋转 45° cube 并集：拓扑有效 \\
  \bottomrule
\end{tabular}%
}
\end{center}

全模块共 \textbf{26 个}单元测试，其中 7 个为共面三角形处理测试，12 个为边界场景
（空网格、完全重合、包含关系、非轴对齐）新增测试。

\section{设计取舍}

\begin{itemize}
  \item \textbf{为何用 corefinement 而非纯射线法？} 旧射线法将三角形作为整体
        保留/丢弃，$A - A$ 残留最多 2 面。Corefinement 在边交点处精确切割，
        使 $A - A$ 严格为空，部分重叠场景几何更精确。
  \item \textbf{为何 7 采样 + 内向偏移？} 重心可能恰在另一网格表面（共面情形），
        单点采样产生歧义。7 采样覆盖面内多点，内向偏移确保共面时采样点
        明确落入网格内部，多数表决鲁棒。
  \item \textbf{为何共面使用 Sutherland-Hodgman 而非更通用的 2D 布尔？}
        Sutherland-Hodgman 仅支持\textbf{凸裁剪多边形}，对三角形 $B$ 而言天然满足。
        相比 Greiner-Hormann 等通用算法，它实现简洁、数值稳定，且仅需 3 次裁剪即可
        得到 $A \cap B$ 的凸多边形。配合内部约束点的扇形分裂，能精确对齐两侧网格的
        共面重叠区域边界，使背对背共面（如小立方体贴大立方体顶面）的并集/差集产生
        流形结果。
  \item \textbf{为何需要共面回退？} 7 采样 + 内向偏移在\textbf{背对背共面}时失效：
        $A$ 的面法向 $+z$、$B$ 的面法向 $-z$，内向偏移使采样点分别向 $\pm z$ 离开对方体积，
        全部判为外部 $\Rightarrow$ 两侧共面面片均被保留 $\Rightarrow$ 非流形。回退用 2D 包含
        检测将共面重叠区域正确判为 inside，避免此问题。
  \item \textbf{为何 \texttt{add\_triangle} 的错误被忽略？} 布尔结果中可能出现
        退化三角形或重复面，\texttt{add\_triangle} 会返回 \texttt{Err}。
        当前实现选择丢弃错误三角形继续，保证调用方拿到非空结果。
  \item \textbf{为何顶点量化去重？} 浮点位置直接比较不可靠（$10^{-16}$ 误差）。
        量化到 $10^{-9}$ 网格既保证共享顶点被识别，又允许微小数值漂移。
\end{itemize}

\end{document}
